% Generated by empirical_application_empluk/run_manuscript_application.R.
% Do not edit numerical entries by hand.

\begin{table}[H]
\centering
\caption{Linear coefficients in the EmplUK application}
\label{tab:empluk_coefficients}
\begin{tabular}{lccc}
\toprule
Covariate & Linear FE & Partially linear FE & Partially linear FD \\
\midrule
Log wages & -0.297 & -0.341 & -0.430 \\
 & (0.126) & (0.129) & (0.163) \\
Log capital & 0.548 & 0.545 & 0.397 \\
 & (0.051) & (0.050) & (0.055) \\
Log output & 0.265 & --- & --- \\
 & (0.153) &  &  \\
\bottomrule
\end{tabular}
\begin{minipage}{0.92\textwidth}
\tablenote\ The dependent variable is log employment in the levels specifications and its first difference in the FD specification. All specifications include year effects; the levels models also include firm effects. Standard errors clustered by firm are reported in parentheses. In the partially linear models, log output enters through a smooth function and therefore has no scalar coefficient; its estimated EDF is 2.87 in the factor fixed-effects specification and 0.96 in the first-difference specification.
\end{minipage}
\end{table}

\begin{table}[H]
\centering
\footnotesize
\caption{Comparison of specifications in the EmplUK application}
\label{tab:empluk_model_comparison}
\begin{tabular}{lrrrrrrr}
\toprule
Model & Wage EDF & Capital EDF & Output EDF & Model df & RMSE & AIC & BIC \\
\midrule
\multicolumn{8}{l}{\textit{Panel A: levels}} \\
\midrule
Linear FE & 1.00 & 1.00 & 1.00 & 152.00 & 0.1180 & -1177.4 & -426.8 \\
Partially linear FE & 1.00 & 1.00 & 2.87 & 154.68 & 0.1172 & -1185.5 & -421.6 \\
Additive GAM FE & 8.84 & 11.61 & 2.51 & 175.68 & 0.1111 & -1252.9 & -385.3 \\
\midrule
\multicolumn{8}{l}{\textit{Panel B: first differences}} \\
\midrule
Partially linear FD & 1.00 & 1.00 & 0.96 & 12.00 & 0.1074 & -1424.1 & -1366.6 \\
\bottomrule
\end{tabular}
\begin{minipage}{0.94\textwidth}
\tablenote\ Linear components are assigned one degree of freedom. Model df is the effective parameter count returned by \code{logLik(fit)} and used in AIC and BIC; when available, it uses \code{mgcv}'s corrected EDF and therefore need not equal a simple sum of the displayed component EDFs. It includes year effects, the estimated scale parameter, and firm effects in the levels models. RMSE is calculated for log employment in Panel A and its first difference in Panel B. AIC and BIC are reported descriptively; fREML selects the smoothing parameters. RMSE, AIC, and BIC should not be compared across panels because the dependent variables, estimation samples, and objectives differ.
\end{minipage}
\end{table}
